% ══════════════════════════════════════════════════════════════════
%  FICHE D'EXERCICES — Chapitre 7 : Les acides carboxyliques et leurs dérivés
% ══════════════════════════════════════════════════════════════════
\section{Exercices type Baccalauréat}
\textit{Données : $M(\ce{H})=1$, $M(\ce{C})=12$, $M(\ce{O})=16$, $M(\ce{Cl})=35{,}5$
\si{\gram\per\mol}.}

\rubrique{Nomenclature, acidité}

\exo{}\diff{1} Nommer : a)~$\ce{HCOOH}$ ; b)~$\ce{CH3-COOH}$ ;
c)~$\ce{CH3-CH2-COOH}$ ; d)~$\ce{(CH3)2CH-COOH}$.

\exo{}\diff{1} Écrire l'équation de la réaction de l'acide éthanoïque avec :
a)~l'hydroxyde de sodium ; b)~le carbonate de sodium $\ce{Na2CO3}$. Nommer le sel formé.

\exo{}\diff{2} Une solution d'acide méthanoïque (formique) de concentration
$C = \SI{0.1}{\mol\per\litre}$ a un pH de $2{,}4$. Montrer que cet acide est faible.

\rubrique{Estérification \& hydrolyse}

\exo{}\diff{1} Écrire l'équation de l'estérification entre l'acide éthanoïque et le
méthanol. Nommer l'ester. Quelle est la réaction inverse ?

\exo{}\diff{2} On mélange \SI{1}{\mol} d'acide éthanoïque et \SI{1}{\mol} d'éthanol.
À l'équilibre, il se forme \SI{0.67}{\mol} d'ester.
\begin{enumerate}[label=\arabic*)]
\item Écrire l'équation de la réaction et dresser le tableau d'avancement.
\item Calculer le rendement de l'estérification.
\item Calculer la constante d'équilibre $K$ de la réaction.
\end{enumerate}

\exo{}\diff{2} L'hydrolyse d'un ester $E$ de masse molaire $M = \SI{88}{\gram\per\mol}$
donne un acide $A$ et un alcool $B$. $A$ est l'acide éthanoïque.
\begin{enumerate}[label=\arabic*)]
\item Déterminer la formule et le nom de $B$.
\item Identifier $E$, le nommer, et écrire l'équation de son hydrolyse.
\end{enumerate}

\rubrique{Dérivés d'acides}

\exo{}\diff{2} Pour préparer un ester avec un meilleur rendement qu'avec l'acide, on
utilise le chlorure d'éthanoyle $\ce{CH3-COCl}$ ou l'anhydride éthanoïque.
\begin{enumerate}[label=\arabic*)]
\item Écrire l'équation de la réaction du chlorure d'éthanoyle avec l'éthanol.
\item Écrire l'équation de la réaction de l'anhydride éthanoïque avec l'éthanol.
\item Citer un avantage de ces réactions par rapport à l'estérification directe.
\end{enumerate}

\exo{}\diff{2} Le chlorure d'éthanoyle réagit avec l'ammoniac $\ce{NH3}$ pour donner
un amide. Écrire l'équation et nommer le produit (éthanamide).

\rubrique{Saponification}

\exo{}\diff{2} On chauffe un ester $E$ avec une solution d'hydroxyde de sodium.
\begin{enumerate}[label=\arabic*)]
\item Comment appelle-t-on cette réaction ? Est-elle limitée ? Justifier.
\item Écrire l'équation de la saponification de l'éthanoate d'éthyle par la soude.
\end{enumerate}

\rubrique{Problèmes de synthèse — type Baccalauréat}

\sujetbac[d'après Baccalauréat D, Cameroun]
Un acide carboxylique saturé $A$ a une masse molaire $M = \SI{74}{\gram\per\mol}$.
\begin{enumerate}[label=\arabic*)]
\item Déterminer la formule brute de $A$ (formule générale $\ce{C_nH_{2n}O2}$), sa
formule semi-développée et son nom.
\item On fait réagir $A$ avec le propan-1-ol. Écrire l'équation de l'estérification
et nommer l'ester $E$ obtenu.
\item Pour améliorer le rendement, on remplace $A$ par son chlorure d'acyle. Donner
la formule et le nom de ce chlorure d'acyle, puis écrire l'équation de sa réaction
avec le propan-1-ol.
\item On veut obtenir \SI{13.2}{\gram} de l'ester $E$. Quelle quantité de matière de
chlorure d'acyle faut-il, en supposant la réaction totale ?
\end{enumerate}

\sujetbac[Détermination par dosage]
On dose \SI{20}{\milli\litre} d'une solution d'un monoacide carboxylique $A$ par une
solution d'hydroxyde de sodium de concentration $C_b = \SI{0.10}{\mol\per\litre}$.
L'équivalence est atteinte pour un volume versé $V_{bE} = \SI{16}{\milli\litre}$.
La solution d'acide a été préparée en dissolvant \SI{0.192}{\gram} de $A$ dans
\SI{20}{\milli\litre} d'eau.
\begin{enumerate}[label=\arabic*)]
\item Écrire l'équation de la réaction de dosage.
\item Déterminer la quantité de matière, puis la concentration de l'acide $A$.
\item En déduire la masse molaire de $A$, sa formule et son nom.
\end{enumerate}
